%% =================================================================
%% Kingma, Ba. "Adam: A Method for Stochastic Optimization."
%% ICLR 2015. arXiv:1412.6980v9.
%%
%% Reconstruction of page 6 (printed p. 6) as study notes.
%% Generated by: reproduce-all-pages-basic-tex (batch 1-15)
%%
%% Structure: Figure 1 (two-panel training-cost comparison) + prose
%% + Sec 6.2 EXPERIMENT: MULTI-LAYER NEURAL NETWORKS + Sec 6.3
%% EXPERIMENT: CONVOLUTIONAL NEURAL NETWORKS (cut off).
%% No equations on this page.
%% =================================================================

\documentclass[11pt]{article}

\usepackage[margin=1in]{geometry}
\usepackage{amsmath, amssymb}
\usepackage{mathrsfs}
\usepackage{bm}
\usepackage{xcolor}
\usepackage{fancyhdr}

\pagestyle{fancy}
\fancyhf{}
\lhead{\small Published as a conference paper at ICLR 2015}
\rhead{}
\cfoot{\small 6 $\to$ \arabic{page}}
\renewcommand{\headrulewidth}{0.4pt}

\setcounter{page}{1}
\setcounter{equation}{5}

\begin{document}

% ---------- Figure 1 placeholder ----------
\noindent
\begin{center}
\fbox{\parbox{0.92\textwidth}{\centering\vspace{1em}
\textbf{[Figure 1 --- placeholder]}\\[0.8em]
Two side-by-side training-cost curves for L2-regularized logistic
regression, showing Adam vs.\ baseline optimizers.\\[0.4em]
\textbf{Left panel --- MNIST Logistic Regression}: training cost
(y-axis, $\approx 0.2$ to $0.7$) vs.\ iterations over entire
dataset (x-axis, $0$ to $\sim 45$). Three curves: \emph{AdaGrad}
(slowest), \emph{SGDNesterov}, and \emph{Adam} (both SGDNesterov
and Adam converge faster than AdaGrad; Adam and SGDNesterov are
visually comparable).\\[0.3em]
\textbf{Right panel --- IMDB BoW feature Logistic Regression}:
training cost (y-axis, $\approx 0.25$ to $0.50$) vs.\ iterations
over entire dataset (x-axis, $0$ to $\sim 160$). Four curves, all
with dropout applied to the sparse BoW features: \emph{Adagrad +
dropout}, \emph{RMSProp + dropout}, \emph{SGDNesterov + dropout},
and \emph{Adam + dropout}. The RMSProp curve is visibly noisier
(larger fluctuations) than the others; Adam + dropout converges
to the lowest training cost overall, with Adagrad+dropout
tracking closely.
\vspace{1em}}}
\end{center}
\vspace{0.3em}
\begin{center}
\textbf{Figure 1:} Logistic regression training negative log
likelihood on MNIST images and IMDB movie reviews with 10,000
bag-of-words (BoW) feature vectors.
\end{center}

\vspace{1em}

% ---------- Prose after Fig. 1 ----------
\noindent
training to prevent over-fitting. In figure~1, Adagrad outperforms
SGD with Nesterov momentum by a large margin both with and
without dropout noise. Adam converges as fast as Adagrad. The
empirical performance of Adam is consistent with our theoretical
findings in sections~2 and~4. Similar to Adagrad, Adam can take
advantage of sparse features and obtain faster convergence rate
than normal SGD with momentum.

\vspace{0.5em}

\subsection*{6.2 \textsc{Experiment: Multi-layer neural networks}}

Multi-layer neural network are powerful models with non-convex
objective functions. Although our convergence analysis does not
apply to non-convex problems, we empirically found that Adam
often outperforms other methods in such cases. In our
experiments, we made model choices that are consistent with
previous publications in the area; a neural network model with
two fully connected hidden layers with $1000$ hidden units each
and ReLU activation are used for this experiment with minibatch
size of $128$.

First, we study different optimizers using the standard
deterministic cross-entropy objective function with $L_2$ weight
decay on the parameters to prevent over-fitting. The
sum-of-functions (SFO) method (Sohl-Dickstein et al., 2014) is a
recently proposed quasi-Newton method that works with minibatches
of data and has shown good performance on optimization of
multi-layer neural networks. We used their implementation and
compared with Adam to train such models. Figure~2 shows that
Adam makes faster progress in terms of both the number of
iterations and wall-clock time. Due to the cost of updating
curvature information, SFO is 5--10x slower per iteration
compared to Adam, and has a memory requirement that is linear in
the number minibatches.

Stochastic regularization methods, such as dropout, are an
effective way to prevent over-fitting and often used in practice
due to their simplicity. SFO assumes deterministic
subfunctions, and indeed failed to converge on cost functions
with stochastic regularization. We compare the effectiveness of
Adam to other stochastic first order methods on multi-layer
neural networks trained with dropout noise. Figure~2 shows our
results; Adam shows better convergence than other methods.

\vspace{0.5em}

\subsection*{6.3 \textsc{Experiment: Convolutional neural networks}}

Convolutional neural networks (CNNs) with several layers of
convolution, pooling and non-linear units have shown considerable
success in computer vision tasks. Unlike most fully connected
neural nets, weight sharing in CNNs results in vastly different
gradients in different layers. A smaller learning rate for the
convolution layers is often used in practice when applying SGD.
We show the effectiveness of Adam in deep CNNs. Our CNN
architecture has three alternating stages of $5 \times 5$
convolution filters and $3 \times 3$ max pooling with stride of
$2$ that are followed by a fully connected layer of $1000$
rectified linear hidden units (ReLU's). The input image are
pre-processed by whitening, and

% [page ends here — continues on page 7]

\end{document}
